% Source: https://artofproblemsolving.com/community/c3h1269976_area_confusion
% Stable identity: aops-topic-1269976-q01
Let the side containing the wall have length $x$ and the other side $y$.

\emph{Case $x\ge 20$.} The whole wall is used, and the fence covers the remaining $x-20$ of that side, the opposite side $x$, and the two sides $y$: $(x-20)+x+2y=100$, so $x+y=60$. The area $xy$ with $x+y=60$ is largest at $x=y=30$, giving $900$.

\emph{Case $x<20$.} Only part of the wall is used and $x+2y=100$, so the area is $x\cdot\dfrac{100-x}{2}$, which increases for $x<50$; its value at $x=20$ is $800$.

The maximum is $\boxed{900}$ square feet (a $30\times 30$ square with the whole wall in one side). The poster's $625$ is the best square made of fence alone, which ignores the wall.

\begin{asy}
usepackage("amsmath");
size(2.8cm);
defaultpen(fontsize(10pt));
draw((0,0)--(3,0)--(3,3)--(0,3)--cycle,blue+linewidth(1)); draw((0.5,3)--(2.5,3),linewidth(3)); label("wall $20$",(1.5,3),N); label("$x$",(1.5,0),S); label("$y$",(3,1.5),E);
label("\mbox{}",(min(currentpicture,user=true)+max(currentpicture,user=true))/2,black); // keeps the text after the figure black when LaTeX typesets the labels inline
\end{asy}
